\documentclass[a4paper,landscape]{article}
\usepackage[utf8]{inputenc}
\usepackage{listings}
\usepackage{multicol}
\usepackage{fancyhdr}
\usepackage{geometry}
\usepackage{graphicx}
\usepackage{tocloft}
\usepackage{tikz}
\geometry{a4paper, landscape, margin=0.7in}
\pagestyle{empty}


\usepackage{xcolor}
% Define custom colors
\definecolor{codebg}{RGB}{248,248,248}
\definecolor{keywordcolor}{RGB}{0,0,180}
\definecolor{stringcolor}{RGB}{163,21,21}
\definecolor{commentcolor}{RGB}{0,128,0}
\definecolor{numbercolor}{RGB}{120,120,120}
\definecolor{identifiercolor}{RGB}{0,0,0}

\lstset{
    language=C++,
    backgroundcolor=\color{codebg},
    basicstyle=\footnotesize\ttfamily\color{identifiercolor},
    keywordstyle=\color{keywordcolor}\bfseries,
    stringstyle=\color{stringcolor},
    commentstyle=\color{commentcolor}\itshape,
    numberstyle=\tiny\color{numbercolor},
    numbers=left,
    stepnumber=1,
    numbersep=8pt,
    breaklines=true,
    showstringspaces=false,
    columns=fullflexible,
    tabsize=2,
    rulecolor=\color{black},
    captionpos=b
}

\setlength{\columnsep}{1cm}  % Set space between columns
\columnseprule=0.4pt  % Add vertical line with thickness 0.4pt
\begin{document}
\pagestyle{fancy}
\fancyhf{} % clear all header and footer
\fancyfoot[C]{\thepage} % center page number in footer
\renewcommand{\headrulewidth}{0pt} % remove header line
\renewcommand{\footrulewidth}{0pt} % remove footer line (optional)
\begin{titlepage}
    \centering
    \vspace*{3cm}
    {\Huge\bfseries Personal Handbook\par}
    \vspace{1.5cm}
    \begin{multicols}{3}
    \renewcommand{\contentsname}{}
    \tableofcontents
    \addtocontents{toc}{\protect\thispagestyle{empty}}  % Optional: remove page number
    \end{multicols}
\end{titlepage}
\newpage
\begin{multicols}{2}
\hline
\section{Bridge}
\hline
\begin{lstlisting}[language=C++]
void IS_BRIDGE(int v,int to); // some function to process the found bridge. Printing for instance
int n; // number of nodes
vector<vector<int>> adj; // adjacency list of graph

vector<bool> visited;
vector<int> tin, low;
int timer;

void dfs(int v, int p = -1) {
    visited[v] = true;
    tin[v] = low[v] = timer++;
    bool parent_skipped = false;
    for (int to : adj[v]) {
        if (to == p && !parent_skipped) {
            parent_skipped = true;
            continue;
        }
        if (visited[to]) {
            low[v] = min(low[v], tin[to]);
        } else {
            dfs(to, v);
            low[v] = min(low[v], low[to]);
            if (low[to] > tin[v])
                IS_BRIDGE(v, to);
        }
    }
}

void find_bridges() {
    timer = 0;
    visited.assign(n, false);
    tin.assign(n, -1);
    low.assign(n, -1);
    for (int i = 0; i < n; ++i) {
        if (!visited[i])
            dfs(i);
    }
}
\end{lstlisting}
\hline
\section{Articulation point}
\hline
\begin{lstlisting}[language=C++]
void IS_CUTPOINT(int v); // here v is articulation point, i can print for instance.
int n; // number of nodes
vector<vector<int>> adj; // adjacency list of graph

vector<bool> visited;
vector<int> tin, low;
int timer;

void dfs(int v, int p = -1) {
    visited[v] = true;
    tin[v] = low[v] = timer++;
    int children=0;
    for (int to : adj[v]) {
        if (to == p) continue;
        if (visited[to]) {
            low[v] = min(low[v], tin[to]);
        } else {
            dfs(to, v);
            low[v] = min(low[v], low[to]);
            if (low[to] >= tin[v] && p!=-1)
                IS_CUTPOINT(v);
            ++children;
        }
    }
    if(p == -1 && children > 1)
        IS_CUTPOINT(v);
}

void find_cutpoints() {
    timer = 0;
    visited.assign(n, false);
    tin.assign(n, -1);
    low.assign(n, -1);
    for (int i = 0; i < n; ++i) {
        if (!visited[i])
            dfs (i);
    }
}
\end{lstlisting}
\hline
\section{Line container}
\hline
\begin{lstlisting}[language=C++]
struct Line {
	mutable int k, m, p;
	bool operator<(const Line& o) const { return k < o.k; }
	bool operator<(int x) const { return p < x; }
};
 
// queries for maximum
// if we want minimum then make minimum = true
bool minimum = true;
struct LineContainer : multiset<Line, less<>> {
	// (for doubles, use inf = 1/.0, div(a,b) = a/b)
	static const int inf = LONG_MAX;
	int div(int a, int b) { // floored division
		return a / b - ((a ^ b) < 0 && a % b); }
	bool isect(iterator x, iterator y) {
		if (y == end()) return x->p = inf, 0;
		if (x->k == y->k) x->p = x->m > y->m ? inf : -inf;
		else x->p = div(y->m - x->m, x->k - y->k);
		return x->p >= y->p;
	}
	// y = m*x+c; x wiint be the query parameter used later
	void add(int m, int c) {
	    if( minimum ){
            m=-m;
            c=-c;
	    }
		auto z = insert({m, c, 0}), y = z++, x = y;
		while (isect(y, z)) z = erase(z);
		if (x != begin() && isect(--x, y)) isect(x, y = erase(y));
		while ((y = x) != begin() && (--x)->p >= y->p)
			isect(x, erase(y));
	}
	int get(int x) {
		assert(!empty());
		auto l = *lower_bound(x);
		if( minimum ) return -(l.k * x + l.m);
		return l.k * x + l.m;
	}
};
\end{lstlisting}
\hline
\section{Convex Hull}
\hline
\begin{lstlisting}
struct pt {
    double x, y;
    bool operator == (pt const& t) const {
        return x == t.x && y == t.y;
    }
};

int orientation(pt a, pt b, pt c) {
    double v = a.x*(b.y-c.y)+b.x*(c.y-a.y)+c.x*(a.y-b.y);
    if (v < 0) return -1; // clockwise
    if (v > 0) return +1; // counter-clockwise
    return 0;
}

bool cw(pt a, pt b, pt c, bool include_collinear) {
    int o = orientation(a, b, c);
    return o < 0 || (include_collinear && o == 0);
}
bool collinear(pt a, pt b, pt c) { return orientation(a, b, c) == 0; }

void convex_hull(vector<pt>& a, bool include_collinear = false) {
    pt p0 = *min_element(a.begin(), a.end(), [](pt a, pt b) {
        return make_pair(a.y, a.x) < make_pair(b.y, b.x);
    });
    sort(a.begin(), a.end(), [&p0](const pt& a, const pt& b) {
        int o = orientation(p0, a, b);
        if (o == 0)
            return (p0.x-a.x)*(p0.x-a.x) + (p0.y-a.y)*(p0.y-a.y)
                < (p0.x-b.x)*(p0.x-b.x) + (p0.y-b.y)*(p0.y-b.y);
        return o < 0;
    });
    if (include_collinear) {
        int i = (int)a.size()-1;
        while (i >= 0 && collinear(p0, a[i], a.back())) i--;
        reverse(a.begin()+i+1, a.end());
    }

    vector<pt> st;
    for (int i = 0; i < (int)a.size(); i++) {
        while (st.size() > 1 && !cw(st[st.size()-2], st.back(), a[i], include_collinear))
            st.pop_back();
        st.push_back(a[i]);
    }

    if (include_collinear == false && st.size() == 2 && st[0] == st[1])
        st.pop_back();

    a = st;
}
\end{lstlisting}
\hline
\section{KMP}
\hline
\begin{lstlisting}
vector<int> prefix_function(string s) {
    int n = (int)s.length();
    vector<int> pi(n);
    for (int i = 1; i < n; i++) {
        int j = pi[i-1];
        while (j > 0 && s[i] != s[j])
            j = pi[j-1];
        if (s[i] == s[j])
            j++;
        pi[i] = j;
    }
    return pi;
}
\end{lstlisting}

\hline
\section{Z function}
\hline
\begin{lstlisting}
vector<int> z_function(string s) {
    int n = s.size();
    vector<int> z(n);
    int l = 0, r = 0;
    for(int i = 1; i < n; i++) {
        if(i < r) {
            z[i] = min(r - i, z[i - l]);
        }
        while(i + z[i] < n && s[z[i]] == s[i + z[i]]) {
            z[i]++;
        }
        if(i + z[i] > r) {
            l = i;
            r = i + z[i];
        }
    }
    return z;
}
\end{lstlisting}
\hline
\section{Suffix Array}
\hline
\begin{lstlisting}
vector<int> sort_cyclic_shifts(string const& s) {
    int n = s.size();
    const int alphabet = 256;
    vector<int> p(n), c(n), cnt(max(alphabet, n), 0);
    for (int i = 0; i < n; i++)
        cnt[s[i]]++;
    for (int i = 1; i < alphabet; i++)
        cnt[i] += cnt[i-1];
    for (int i = 0; i < n; i++)
        p[--cnt[s[i]]] = i;
    c[p[0]] = 0;
    int classes = 1;
    for (int i = 1; i < n; i++) {
        if (s[p[i]] != s[p[i-1]])
            classes++;
        c[p[i]] = classes - 1;
    }
    vector<int> pn(n), cn(n);
    for (int h = 0; (1 << h) < n; ++h) {
        for (int i = 0; i < n; i++) {
            pn[i] = p[i] - (1 << h);
            if (pn[i] < 0)
                pn[i] += n;
        }
        fill(cnt.begin(), cnt.begin() + classes, 0);
        for (int i = 0; i < n; i++)
            cnt[c[pn[i]]]++;
        for (int i = 1; i < classes; i++)
            cnt[i] += cnt[i-1];
        for (int i = n-1; i >= 0; i--)
            p[--cnt[c[pn[i]]]] = pn[i];
        cn[p[0]] = 0;
        classes = 1;
        for (int i = 1; i < n; i++) {
            pair<int, int> cur = {c[p[i]], c[(p[i] + (1 << h)) % n]};
            pair<int, int> prev = {c[p[i-1]], c[(p[i-1] + (1 << h)) % n]};
            if (cur != prev)
                ++classes;
            cn[p[i]] = classes - 1;
        }
        c.swap(cn);
    }
    return p;
}
\end{lstlisting}
\hline
\section{Manacher}
\hline 
\begin{lstlisting}
vector<int> manacher_odd(string s) {
    int n = s.size();
    s = "$" + s + "^";
    vector<int> p(n + 2);
    int l = 0, r = 1;
    for(int i = 1; i <= n; i++) {
        p[i] = min(r - i, p[l + (r - i)]);
        while(s[i - p[i]] == s[i + p[i]]) {
            p[i]++;
        }
        if(i + p[i] > r) {
            l = i - p[i], r = i + p[i];
        }
    }
    return vector<int>(begin(p) + 1, end(p) - 1);
}
\end{lstlisting}
\hline
\section{Short ST}
\hline
\begin{lstlisting}
const int N = 1e5;  // limit for array size
int n;  // array size
int t[2 * N];

void build() {  // build the tree (first set positions n through 2n-1 as original array
  for (int i = n - 1; i > 0; --i) t[i] = t[i<<1] + t[i<<1|1];
}

void modify(int p, int value) {  // set value at position p
  for (t[p += n] = value; p > 1; p >>= 1) t[p>>1] = t[p] + t[p^1];
}

int query(int l, int r) {  // sum on interval [l, r)
  int res = 0;
  for (l += n, r += n; l < r; l >>= 1, r >>= 1) {
    if (l&1) res += t[l++];
    if (r&1) res += t[--r];
  }
  return res;
}
\end{lstlisting}
\hline
\section{Lazy ST}
\hline
\begin{lstlisting}
/**
 * ===================================================
 * USABILITY GUIDE
 * ===================================================
 * 1. Define Combine Operation:
 *    - Sum: combine(a,b) = a + b
 *    - Max: combine(a,b) = max(a,b)
 *    - Min: combine(a,b) = min(a,b)
 *    - Xor: combine(a,b) = a ^ b
 *
 * 2. Define Lazy Meaning:
 *    - Range Assignment: lazy[node] = v
 *    - Range Add: lazy[node] += v
 *    - Range Max-Update (chmax): lazy[node] = max(lazy[node], v)
 *
 * 3. Modify update_lazy():
 *    Example:
 *      - Range Assignment + Sum:
 *          val[node] = lazy[node];
 *          sum[node] = (r-l+1) * lazy[node];
 *
 *      - Range Add + Sum:
 *          val[node] += lazy[node];
 *          sum[node] += (r-l+1) * lazy[node];
 *
 *      - Range Assignment + Max:
 *          val[node] = lazy[node];
 *          sum[node] = lazy[node];  // here sum[] stores max
 *
 * 4. Adjust Query Neutral Value:
 *    - Sum: return 0
 *    - Max: return -INF
 *    - Min: return +INF
 *    - Xor: return 0
 *    - GCD: return 0
 *
 * Complexity: O(log N) per operation.
 * Root is at index 1, responsible for array[0..n-1].
 */
struct segTree {
    int n;
    vector<int> val, sum, lazy;

    segTree(int nn) {
        n = nn;
        val.resize(n << 2);
        sum.resize(n << 2);
        lazy.resize(n << 2);
    }

    void build(vector<int>& a, int node, int l, int r) {
        if (l == r) {
            val[node] = a[l];
            sum[node] = a[l];   // adjust if using other ops
        } else {
            int m = (l + r) >> 1;
            build(a, node << 1, l, m);
            build(a, (node << 1) | 1, m + 1, r);
            val[node] = val[node << 1] + val[(node << 1) | 1]; // combine()
            sum[node] = sum[node << 1] + sum[(node << 1) | 1]; // combine()
        }
    }

    void update_lazy(int node, int l, int r) {
        if (!lazy[node]) return;
        int v = lazy[node];

        // Example: Range assignment with sum
        val[node] = v;
        sum[node] = (r - l + 1) * v;

        // For other ops: replace this block
    }

    void push(int node, int l, int r) {
        if (lazy[node]) {
            lazy[node << 1] = lazy[(node << 1) | 1] = lazy[node];
            update_lazy(node, l, r);

            int m = (l + r) >> 1;
            update_lazy(node << 1, l, m);
            update_lazy((node << 1) | 1, m + 1, r);
            lazy[node] = 0;
        }
    }

    void assign_range(int node, int tl, int tr, int l, int r, int v) {
        if (l > r) return;
        if (tl == l && tr == r) {
            lazy[node] = v;
            if (tl == tr) update_lazy(node, tl, tr);
            else push(node, tl, tr);
        } else {
            push(node, tl, tr);
            int tm = (tl + tr) >> 1;
            assign_range(node << 1, tl, tm, l, min(tm, r), v);
            assign_range((node << 1) | 1, tm + 1, tr, max(tm + 1, l), r, v);
            sum[node] = sum[node << 1] + sum[(node << 1) | 1]; // combine()
        }
    }

    int range_query(int node, int tl, int tr, int l, int r) {
        if (l > r) return 0;  // neutral element for sum
        if (tl == l && tr == r) {
            if (tl == tr) update_lazy(node, tl, tr);
            else push(node, tl, tr);
            return sum[node];
        }
        push(node, tl, tr);
        int tm = (tl + tr) >> 1;
        return range_query(node << 1, tl, tm, l, min(r, tm)) +
               range_query((node << 1) | 1, tm + 1, tr, max(tm + 1, l), r); // combine()
    }

    int point_query(int node, int tl, int tr, int pos) {
        if (tl > tr) return 0;
        if (tl == pos && tr == pos) {
            update_lazy(node, tl, tr);
            return val[node];
        }
        push(node, tl, tr);
        int tm = (tl + tr) >> 1;
        if (pos <= tm)
            return point_query(node << 1, tl, tm, pos);
        return point_query((node << 1) | 1, tm + 1, tr, pos);
    }
};

\end{lstlisting}
\hline
\section{Sieve}
\hline
\begin{lstlisting}
const int MAX_PR = 5'000'000;
bitset<MAX_PR> isprime;
vi eratosthenesSieve(int lim) {
	isprime.set(); isprime[0] = isprime[1] = 0;
	for (int i = 4; i < lim; i += 2) isprime[i] = 0;
	for (int i = 3; i*i < lim; i += 2) if (isprime[i])
		for (int j = i*i; j < lim; j += i*2) isprime[j] = 0;
	vi pr;
	rep(i,2,lim) if (isprime[i]) pr.push_back(i);
	return pr;
}
\end{lstlisting}
\hline
\section{Fraction finder}
\hline
\begin{lstlisting}
/**
 * Description: Given $f$ and $N$, finds the smallest fraction $p/q \in [0, 1]$
 * such that $f(p/q)$ is true, and $p, q \le N$.
 * You may want to throw an exception from $f$ if it finds an exact solution,
 * in which case $N$ can be removed.
 * Usage: fracBS([](Frac f) { return f.p>=3*f.q; }, 10); // {1,3}
 * Time: O(\log(N))
 * Status: stress-tested for n <= 300
 */

struct Frac { ll p, q; };

template<class F>
Frac fracBS(F f, ll N) {
	bool dir = 1, A = 1, B = 1;
	Frac lo{0, 1}, hi{1, 1}; // Set hi to 1/0 to search (0, N]
	if (f(lo)) return lo;
	assert(f(hi));
	while (A || B) {
		ll adv = 0, step = 1; // move hi if dir, else lo
		for (int si = 0; step; (step *= 2) >>= si) {
			adv += step;
			Frac mid{lo.p * adv + hi.p, lo.q * adv + hi.q};
			if (abs(mid.p) > N || mid.q > N || dir == !f(mid)) {
				adv -= step; si = 2;
			}
		}
		hi.p += lo.p * adv;
		hi.q += lo.q * adv;
		dir = !dir;
		swap(lo, hi);
		A = B; B = !!adv;
	}
	return dir ? hi : lo;
}
\end{lstlisting}
\hline
\section{FFT v2}
\hline
\begin{lstlisting}
/**
 * Description: fft(a) computes $\hat f(k) = \sum_x a[x] \exp(2\pi i \cdot k x / N)$ for all $k$. N must be a power of 2.
   Useful for convolution:
   \texttt{conv(a, b) = c}, where $c[x] = \sum a[i]b[x-i]$.
   For convolution of complex numbers or more than two vectors: FFT, multiply
   pointwise, divide by n, reverse(start+1, end), FFT back.
   Rounding is safe if $(\sum a_i^2 + \sum b_i^2)\log_2{N} < 9\cdot10^{14}$
   (in practice $10^{16}$; higher for random inputs).
   Otherwise, use NTT/FFTMod.
 * Time: O(N \log N) with $N = |A|+|B|$ ($\tilde 1s$ for $N=2^{22}$)
 * Status: somewhat tested
 */
#pragma once

typedef complex<double> C;
typedef vector<double> vd;
void fft(vector<C>& a) {
	int n = sz(a), L = 31 - __builtin_clz(n);
	static vector<complex<long double>> R(2, 1);
	static vector<C> rt(2, 1);  // (^ 10% faster if double)
	for (static int k = 2; k < n; k *= 2) {
		R.resize(n); rt.resize(n);
		auto x = polar(1.0L, acos(-1.0L) / k);
		rep(i,k,2*k) rt[i] = R[i] = i&1 ? R[i/2] * x : R[i/2];
	}
	vi rev(n);
	rep(i,0,n) rev[i] = (rev[i / 2] | (i & 1) << L) / 2;
	rep(i,0,n) if (i < rev[i]) swap(a[i], a[rev[i]]);
	for (int k = 1; k < n; k *= 2)
		for (int i = 0; i < n; i += 2 * k) rep(j,0,k) {
			// C z = rt[j+k] * a[i+j+k]; // (25% faster if hand-rolled)  /// include-line
			auto x = (double *)&rt[j+k], y = (double *)&a[i+j+k];        /// exclude-line
			C z(x[0]*y[0] - x[1]*y[1], x[0]*y[1] + x[1]*y[0]);           /// exclude-line
			a[i + j + k] = a[i + j] - z;
			a[i + j] += z;
		}
}
vd conv(const vd& a, const vd& b) {
	if (a.empty() || b.empty()) return {};
	vd res(sz(a) + sz(b) - 1);
	int L = 32 - __builtin_clz(sz(res)), n = 1 << L;
	vector<C> in(n), out(n);
	copy(all(a), begin(in));
	rep(i,0,sz(b)) in[i].imag(b[i]);
	fft(in);
	for (C& x : in) x *= x;
	rep(i,0,n) out[i] = in[-i & (n - 1)] - conj(in[i]);
	fft(out);
	rep(i,0,sz(res)) res[i] = imag(out[i]) / (4 * n);
	return res;
}
\end{lstlisting}
\hline
\section{FFT Mod}
\hline
\begin{lstlisting}
/**
 * Description: Higher precision FFT, can be used for convolutions modulo arbitrary integers
 * as long as $N\log_2N\cdot \text{mod} < 8.6 \cdot 10^{14}$ (in practice $10^{16}$ or higher).
 * Inputs must be in $[0, \text{mod})$.
 * Time: O(N \log N), where $N = |A|+|B|$ (twice as slow as NTT or FFT)
 * Status: stress-tested
 */
#include "FastFourierTransform.h"
typedef vector<ll> vl;
template<int M> vl convMod(const vl &a, const vl &b) {
	if (a.empty() || b.empty()) return {};
	vl res(sz(a) + sz(b) - 1);
	int B=32-__builtin_clz(sz(res)), n=1<<B, cut=int(sqrt(M));
	vector<C> L(n), R(n), outs(n), outl(n);
	rep(i,0,sz(a)) L[i] = C((int)a[i] / cut, (int)a[i] % cut);
	rep(i,0,sz(b)) R[i] = C((int)b[i] / cut, (int)b[i] % cut);
	fft(L), fft(R);
	rep(i,0,n) {
		int j = -i & (n - 1);
		outl[j] = (L[i] + conj(L[j])) * R[i] / (2.0 * n);
		outs[j] = (L[i] - conj(L[j])) * R[i] / (2.0 * n) / 1i;
	}
	fft(outl), fft(outs);
	rep(i,0,sz(res)) {
		ll av = ll(real(outl[i])+.5), cv = ll(imag(outs[i])+.5);
		ll bv = ll(imag(outl[i])+.5) + ll(real(outs[i])+.5);
		res[i] = ((av % M * cut + bv) % M * cut + cv) % M;
	}
	return res;
}
\end{lstlisting}
\hline
\section{Heavy light Decomp}
\hline
\begin{lstlisting}
include "../data-structures/LazySegmentTree.h"

template <bool VALS_EDGES> struct HLD {
	int N, tim = 0;
	vector<vi> adj;
	vi par, siz, rt, pos;
	Node *tree;
	HLD(vector<vi> adj_)
		: N(sz(adj_)), adj(adj_), par(N, -1), siz(N, 1),
		  rt(N),pos(N),tree(new Node(0, N)){ dfsSz(0); dfsHld(0); }
	void dfsSz(int v) {
		if (par[v] != -1) adj[v].erase(find(all(adj[v]), par[v]));
		for (int& u : adj[v]) {
			par[u] = v;
			dfsSz(u);
			siz[v] += siz[u];
			if (siz[u] > siz[adj[v][0]]) swap(u, adj[v][0]);
		}
	}
	void dfsHld(int v) {
		pos[v] = tim++;
		for (int u : adj[v]) {
			rt[u] = (u == adj[v][0] ? rt[v] : u);
			dfsHld(u);
		}
	}
	template <class B> void process(int u, int v, B op) {
		for (; rt[u] != rt[v]; v = par[rt[v]]) {
			if (pos[rt[u]] > pos[rt[v]]) swap(u, v);
			op(pos[rt[v]], pos[v] + 1);
		}
		if (pos[u] > pos[v]) swap(u, v);
		op(pos[u] + VALS_EDGES, pos[v] + 1);
	}
	void modifyPath(int u, int v, int val) {
		process(u, v, [&](int l, int r) { tree->add(l, r, val); });
	}
	int queryPath(int u, int v) { // Modify depending on problem
		int res = -1e9;
		process(u, v, [&](int l, int r) {
				res = max(res, tree->query(l, r));
		});
		return res;
	}
	int querySubtree(int v) { // modifySubtree is similar
		return tree->query(pos[v] + VALS_EDGES, pos[v] + siz[v]);
	}
};
\end{lstlisting}
\hline
\section{Global Min-Cut}
\hline
\begin{lstlisting}
pair<int, vi> globalMinCut(vector<vi> mat) {
	pair<int, vi> best = {INT_MAX, {}};
	int n = sz(mat);
	vector<vi> co(n);
	rep(i,0,n) co[i] = {i};
	rep(ph,1,n) {
		vi w = mat[0];
		size_t s = 0, t = 0;
		rep(it,0,n-ph) { // O(V^2) -> O(E log V) with prio. queue
			w[t] = INT_MIN;
			s = t, t = max_element(all(w)) - w.begin();
			rep(i,0,n) w[i] += mat[t][i];
		}
		best = min(best, {w[t] - mat[t][t], co[t]});
		co[s].insert(co[s].end(), all(co[t]));
		rep(i,0,n) mat[s][i] += mat[t][i];
		rep(i,0,n) mat[i][s] = mat[s][i];
		mat[0][t] = INT_MIN;
	}
	return best;
}
\end{lstlisting}
\hline
\section{Max Clique}
\hline
\begin{lstlisting}
/**
 * Author: chilli, SJTU, Janez Konc
 * Date: 2019-05-10
 * License: GPL3+
 * Source: https://en.wikipedia.org/wiki/MaxCliqueDyn_maximum_clique_algorithm, https://gitlab.com/janezkonc/mcqd/blob/master/mcqd.h
 * Description: Quickly finds a maximum clique of a graph (given as symmetric bitset
 * matrix; self-edges not allowed). Can be used to find a maximum independent
 * set by finding a clique of the complement graph.
 * Time: Runs in about 1s for n=155 and worst case random graphs (p=.90). Runs
 * faster for sparse graphs.
 * Status: stress-tested
 */
typedef vector<bitset<200>> vb;
struct Maxclique {
	double limit=0.025, pk=0;
	struct Vertex { int i, d=0; };
	typedef vector<Vertex> vv;
	vb e;
	vv V;
	vector<vi> C;
	vi qmax, q, S, old;
	void init(vv& r) {
		for (auto& v : r) v.d = 0;
		for (auto& v : r) for (auto j : r) v.d += e[v.i][j.i];
		sort(all(r), [](auto a, auto b) { return a.d > b.d; });
		int mxD = r[0].d;
		rep(i,0,sz(r)) r[i].d = min(i, mxD) + 1;
	}
	void expand(vv& R, int lev = 1) {
		S[lev] += S[lev - 1] - old[lev];
		old[lev] = S[lev - 1];
		while (sz(R)) {
			if (sz(q) + R.back().d <= sz(qmax)) return;
			q.push_back(R.back().i);
			vv T;
			for(auto v:R) if (e[R.back().i][v.i]) T.push_back({v.i});
			if (sz(T)) {
				if (S[lev]++ / ++pk < limit) init(T);
				int j = 0, mxk = 1, mnk = max(sz(qmax) - sz(q) + 1, 1);
				C[1].clear(), C[2].clear();
				for (auto v : T) {
					int k = 1;
					auto f = [&](int i) { return e[v.i][i]; };
					while (any_of(all(C[k]), f)) k++;
					if (k > mxk) mxk = k, C[mxk + 1].clear();
					if (k < mnk) T[j++].i = v.i;
					C[k].push_back(v.i);
				}
				if (j > 0) T[j - 1].d = 0;
				rep(k,mnk,mxk + 1) for (int i : C[k])
					T[j].i = i, T[j++].d = k;
				expand(T, lev + 1);
			} else if (sz(q) > sz(qmax)) qmax = q;
			q.pop_back(), R.pop_back();
		}
	}
	vi maxClique() { init(V), expand(V); return qmax; }
	Maxclique(vb conn) : e(conn), C(sz(e)+1), S(sz(C)), old(S) {
		rep(i,0,sz(e)) V.push_back({i});
	}
};
\end{lstlisting}
\hline
\section{Dinic 2}
\hline
\begin{lstlisting}
// * $O(\min(E^{1/2}, V^{2/3})E)$ if $U = 1$; $O(\sqrt{V}E)$ for bipartite matching.

struct Dinic {
	struct Edge {
		int to, rev;
		ll c, oc;
		ll flow() { return max(oc - c, 0LL); } // if you need flows
	};
	vi lvl, ptr, q;
	vector<vector<Edge>> adj;
	Dinic(int n) : lvl(n), ptr(n), q(n), adj(n) {}
	void addEdge(int a, int b, ll c, ll rcap = 0) {
		adj[a].push_back({b, sz(adj[b]), c, c});
		adj[b].push_back({a, sz(adj[a]) - 1, rcap, rcap});
	}
	ll dfs(int v, int t, ll f) {
		if (v == t || !f) return f;
		for (int& i = ptr[v]; i < sz(adj[v]); i++) {
			Edge& e = adj[v][i];
			if (lvl[e.to] == lvl[v] + 1)
				if (ll p = dfs(e.to, t, min(f, e.c))) {
					e.c -= p, adj[e.to][e.rev].c += p;
					return p;
				}
		}
		return 0;
	}
	ll calc(int s, int t) {
		ll flow = 0; q[0] = s;
		rep(L,0,31) do { // 'int L=30' maybe faster for random data
			lvl = ptr = vi(sz(q));
			int qi = 0, qe = lvl[s] = 1;
			while (qi < qe && !lvl[t]) {
				int v = q[qi++];
				for (Edge e : adj[v])
					if (!lvl[e.to] && e.c >> (30 - L))
						q[qe++] = e.to, lvl[e.to] = lvl[v] + 1;
			}
			while (ll p = dfs(s, t, LLONG_MAX)) flow += p;
		} while (lvl[t]);
		return flow;
	}
	bool leftOfMinCut(int a) { return lvl[a] != 0; }
};
\end{lstlisting}
\hline
\section{Virtual Tree}
\hline
\begin{lstlisting}
	#include <bits/stdc++.h>
using namespace std;

using ll = long long;

constexpr int MAX_N = 2e5 + 1;
constexpr int LG = 18;
constexpr int MOD = 998244353;

int n, a[MAX_N];
ll dp[MAX_N], ans;
int tin[MAX_N], tout[MAX_N], d[MAX_N], lift[MAX_N][LG], timer;
vector<int> adj[MAX_N], vadj[MAX_N], at_a[MAX_N];

void dfs(int v, int p) {
	at_a[a[v]].push_back(v);
	tin[v] = timer++;

	lift[v][0] = p;
	for (int i = 1; i < LG; i++) { lift[v][i] = lift[lift[v][i - 1]][i - 1]; }

	for (int u : adj[v]) {
		if (u == p) { continue; }
		dfs(u, v);
	}

	tout[v] = timer++;
}

int is_ancestor(int u, int v) { return tin[u] <= tin[v] && tout[v] <= tout[u]; }

int lca(int u, int v) {
	if (is_ancestor(u, v)) { return u; }
	if (is_ancestor(v, u)) { return v; }

	for (int i = LG - 1; i >= 0; i--) {
		if (!is_ancestor(lift[u][i], v)) { u = lift[u][i]; }
	}
	return lift[u][0];
}

bool sort_tin(const int &a, const int &b) { return tin[a] < tin[b]; }

vector<int> vtree(const vector<int> &key) {
	if (key.empty()) return {};

	vector<int> res = key;
	sort(res.begin(), res.end(), sort_tin);

	for (int i = 1; i < (int)key.size(); i++) {
		res.push_back(lca(key[i - 1], key[i]));
	}

	sort(res.begin(), res.end(), sort_tin);
	res.erase(unique(res.begin(), res.end()), res.end());

	for (int v : res) { vadj[v].clear(); }

	for (int i = 1; i < (int)res.size(); i++) {
		vadj[lca(res[i - 1], res[i])].push_back(res[i]);
	}

	return res;
}

int main() {
	cin >> n;

	for (int i = 0; i < n; i++) { cin >> a[i]; }

	for (int i = 1; i < n; i++) {
		int u, v;
		cin >> u >> v;
		u--;
		v--;
		adj[u].push_back(v);
		adj[v].push_back(u);
	}

	dfs(0, 0);

	for (int col = 1; col <= n; col++) {
		vector<int> vt = vtree(at_a[col]);
		reverse(vt.begin(), vt.end());

		for (int v : vt) {
			dp[v] = 1;
			for (int u : vadj[v]) {
				dp[v] *= (dp[u] + 1);
				dp[v] %= MOD;
			}
			if (a[v] != col) { dp[v]--; }

			ans += dp[v];
			ans %= MOD;
			if (a[v] != col) {
				for (int u : vadj[v]) {
					ans += MOD - dp[u];
					ans %= MOD;
				}
			}
		}
	}

	cout << ans << '\n';
}
\end{lstlisting}
\hline
\section{Kruskal Reconstruction Tree}
\hline
\begin{lstlisting}
int getRoot(int u) {
  if (u == dsu[u]) return u;
  return dsu[u] = getRoot(dsu[u]);
}

void addEdge(int u, int v) {
  u = getRoot(u); v = getRoot(v);

  dsu[n] = n;
  dsu[u] = dsu[v] = n;

  adj[n].push_back(u);
  if (u != v) adj[n].push_back(v);

  ++n;
}

void build() {
  for (int i = 0; i < n; ++i) dsu[i] = i;
  for (int i = 0; i < m; ++i) addEdge(U[i], V[i]);
}
\end{lstlisting}
\hline
\section{DSU on Tree(Sack)}
\hline
\begin{lstlisting}
int cnt[maxn];
void dfs(int v, int p, bool keep){
    int mx = -1, bigChild = -1;
    for(auto u : g[v])
       if(u != p && sz[u] > mx)
          mx = sz[u], bigChild = u;
    for(auto u : g[v])
        if(u != p && u != bigChild)
            dfs(u, v, 0);  // run a dfs on small childs and clear them from cnt
    if(bigChild != -1)
        dfs(bigChild, v, 1);  // bigChild marked as big and not cleared from cnt
    for(auto u : g[v])
	if(u != p && u != bigChild)
	    for(int p = st[u]; p < ft[u]; p++)
		cnt[ col[ ver[p] ] ]++;
    cnt[ col[v] ]++;
    //now cnt[c] is the number of vertices in subtree of vertex v that has color c. You can answer the queries easily.
    if(keep == 0)
        for(int p = st[v]; p < ft[v]; p++)
	    cnt[ col[ ver[p] ] ]--;
\end{lstlisting}
\hline
\section{Maximum XOR Subset}
\hline
\begin{lstlisting}
#include<bits/stdc++.h>
using namespace std;
 
#define endl '\n'
 
int basis[35];
void add (int x) {
    for (int i = 31; i >= 0; --i) {
        if ((x >> i) & 1) {
            if (!basis[i]) {
                basis[i] = x;
                return;
            }
            x ^= basis[i];
        }
    }
}
 
int getMaxXor() {
    int ans = 0;
    for (int i = 31; i >= 0; --i) {
        if ( (ans ^ basis[i]) > ans) {
            ans ^= basis[i];
        }
    }
    return ans;
}
void solve () {
    int n, val;
    cin >> n;
    for (int i = 0; i < n; ++i) {
        cin >> val;
        add(val);
    }
    
    cout << getMaxXor() << endl;
}
 
int main() {
    ios::sync_with_stdio(false);
    cin.tie(nullptr);
 
    solve();
 
    return 0;
}
\end{lstlisting}
\hline
\section{Number of Subset XOR}
\hline
\begin{lstlisting}
#include<bits/stdc++.h>
using namespace std;
 
#define endl '\n'
 
int basis[35];
void add (int x) {
    for (int i = 31; i >= 0; --i) {
        if ((x >> i) & 1) {
            if (!basis[i]) {
                basis[i] = x;
                return;
            }
            x ^= basis[i];
        }
    }
}
 
int getXorCount() {
    int ans = 1;
    for (int i = 31; i >= 0; --i) {
        if ( basis[i] > 0) {
            ans <<= 1;
        }
    }
    return ans;
}
void solve () {
    int n, val;
    cin >> n;
    for (int i = 0; i < n; ++i) {
        cin >> val;
        add(val);
    }
    
    cout << getXorCount() << endl;
}
 
int main() {
    ios::sync_with_stdio(false);
    cin.tie(nullptr);
 
    solve();
 
    return 0;
}
\end{lstlisting}
\hline
\section{Critical Cities(Dominator)}
\hline
\begin{lstlisting}
#include <algorithm>
#include <cassert>
#include <cmath>
#include <cstdint>
#include <functional>
#include <iostream>
#include <vector>

using std::cout;
using std::endl;
using std::vector;

// BeginCodeSnip{Tree Class for rdom}
template <typename T> class Tree {
  private:
	const int root = 0;
	const vector<int> &parents;
	const vector<T> &vals;
	const int log2dist;
	vector<vector<int>> pow2ends;
	vector<vector<T>> pow2mins;
	vector<int> depth;

  public:
	Tree(const vector<int> &parents, const vector<T> &vals)
	    : parents(parents), vals(vals),
	      log2dist(std::ceil(std::log2(parents.size() + 1))),
	      pow2ends(parents.size(), vector<int>(log2dist + 1)),
	      pow2mins(parents.size(), vector<T>(log2dist + 1)) {
		assert(parents[root] == -1);

		vector<vector<int>> children(parents.size());
		for (int n = 0; n < parents.size(); n++) {
			if (parents[n] != -1) { children[parents[n]].push_back(n); }
		}
		depth = vector<int>(parents.size());
		vector<int> frontier{root};
		while (!frontier.empty()) {
			int curr = frontier.back();
			frontier.pop_back();
			for (int n : children[curr]) {
				depth[n] = depth[curr] + 1;
				frontier.push_back(n);
			}
		}

		for (int n = 0; n < parents.size(); n++) {
			pow2mins[n][0] = vals[n];
			pow2ends[n][0] = parents[n];
		}
		for (int p = 1; p <= log2dist; p++) {
			for (int n = 0; n < parents.size(); n++) {
				int halfway = pow2ends[n][p - 1];
				if (halfway == -1) {
					pow2ends[n][p] = -1;
					pow2mins[n][p] = pow2mins[n][p - 1];
				} else {
					pow2ends[n][p] = pow2ends[halfway][p - 1];
					pow2mins[n][p] =
					    std::min(pow2mins[n][p - 1], pow2mins[halfway][p - 1]);
				}
			}
		}
	}

	/**
	 * @return the min value on the path from the ancestor to its descendant,
	 *         not including the ancestor.
	 */
	T min_val(int ancestor, int desc) {
		int dist = depth[desc] - depth[ancestor];
		T ret = vals[desc];
		int at = desc;
		for (int pow = 0; pow <= log2dist; pow++) {
			if ((dist & (1 << pow)) != 0) {
				ret = std::min(ret, pow2mins[at][pow]);
				at = pow2ends[at][pow];
			}
		}
		assert(at == ancestor);
		return ret;
	}
};
// EndCodeSnip

int main() {
	int city_num;
	int flight_num;
	std::cin >> city_num >> flight_num;
	vector<vector<int>> adj(city_num);
	vector<vector<int>> rev_adj(city_num);
	for (int f = 0; f < flight_num; f++) {
		int a, b;
		std::cin >> a >> b;
		adj[--a].push_back(--b);
		rev_adj[b].push_back(a);
	}

	// Variables used for the Euler Tour
	vector<int> visit_order;
	// visit_time is initialized with INF so nodes that aren't reachable
	// during the DFS won't mess up our sdom calculation
	vector<int> visit_time(city_num, INT32_MAX);
	vector<int> visit_parent(city_num, -1);
	vector<bool> visited(city_num, false);
	std::function<void(int, int)> dfs;
	dfs = [&](int at, int prev) {
		if (visited[at]) { return; }
		visited[at] = true;

		visit_time[at] = visit_order.size();
		visit_order.push_back(at);
		visit_parent[at] = prev;
		for (int next : adj[at]) { dfs(next, at); }
	};
	dfs(0, -1);

	// We use a function-based interface instead of a class-based one due to
	// heavy reliance on previously calculated values
	vector<int> dsu_parent(city_num, -1);
	vector<int> dsu_min(city_num, INT32_MAX);
	std::function<int(int)> find;
	find = [&](int x) {
		if (dsu_parent[x] == -1) { return x; }
		if (dsu_parent[dsu_parent[x]] != -1) {
			int parent_res = find(dsu_parent[x]);
			if (visit_time[parent_res] < visit_time[dsu_min[x]]) {
				dsu_min[x] = parent_res;
			}
			dsu_parent[x] = dsu_parent[dsu_parent[x]];
		}
		assert(dsu_min[x] != INT32_MAX);
		return dsu_min[x];
	};

	vector<int> sdom(city_num, -1);
	for (int i = visit_order.size() - 1; i > 0; i--) {
		int c = visit_order[i];
		// Iterate over all nodes with a directed edge to c
		for (int from : rev_adj[c]) {
			int find_res = find(from);
			// Take the node with the earliest visit time in the traversal
			if (sdom[c] == -1 || visit_time[find_res] < visit_time[sdom[c]]) {
				sdom[c] = find_res;
			}
		}

		// Link c to its parent in the DFS tree
		dsu_parent[c] = visit_parent[c];
		dsu_min[c] = sdom[c];
	}

	// Initialize the values in the tree
	vector<std::pair<int, int>> sdom_times(city_num, {INT32_MAX, -1});
	for (int c : visit_order) {
		if (c != 0) {
			// The tree takes the minimum of these pairs, but what we
			// really want is the node itself- therefore, we store both
			sdom_times[c] = {visit_time[sdom[c]], c};
		}
	}
	Tree tree(visit_parent, sdom_times);

	vector<int> rdom(city_num, -1);
	for (int c : visit_order) {
		if (c != 0) {
			// Use the definition of the rdom
			rdom[c] = tree.min_val(sdom[c], c).second;
		}
	}

	// Use the properties previously given to compute the idom
	vector<int> idom(city_num, -1);
	for (int c : visit_order) {
		if (c != 0) { idom[c] = rdom[c] == c ? sdom[c] : idom[rdom[c]]; }
	}

	// Trace the idom tree from the finish node to the start node,
	// finding all dominators along the way
	vector<int> critical{0};
	int at = city_num - 1;
	while (at != 0) {
		critical.push_back(at);
		at = idom[at];
	}

	std::sort(critical.begin(), critical.end());

	cout << critical.size() << '\n';
	for (int i = 0; i < critical.size() - 1; i++) { cout << critical[i] + 1 << ' '; }
	cout << critical.back() + 1 << endl;
}
\end{lstlisting}
\hline
\section{DC DP}
\hline
Divide and Conquer is a dynamic programming optimization.

\subsection*{Preconditions}

Some dynamic programming problems have a recurrence of the form:

\[
dp(i, j) = \min_{0 \le k \le j}
\left\{ dp(i - 1, k - 1) + C(k, j) \right\}
\]

where \( C(k, j) \) is a cost function and  
\( dp(i, j) = 0 \) when \( j < 0 \).

Assume:
\[
0 \le i < m, \quad 0 \le j < n
\]
and evaluating \( C \) takes \( O(1) \) time.

A straightforward evaluation runs in:
\[
O(m n^2)
\]
There are \( m \times n \) states and \( n \) transitions per state.

Let \( opt(i, j) \) be the value of \( k \) that minimizes the recurrence.

If the cost function satisfies the \textbf{quadrangle inequality}, then we can prove:

\[
opt(i, j) \le opt(i, j+1)
\]

This is called the \textbf{monotonicity condition}.  
For fixed \( i \), the optimal splitting point increases as \( j \) increases.

This allows us to reduce the search range when computing transitions.

If we know \( opt(i, j) \), then for any \( j' < j \):

\[
opt(i, j') \le opt(i, j)
\]

So when computing \( opt(i, j') \), we only need to search a smaller interval.

\subsection*{Divide and Conquer Optimization}

We apply a divide and conquer strategy:

1. Compute \( opt(i, \lfloor n/2 \rfloor) \).
2. Recursively compute:
   - Left half using upper bound \( opt(i, n/2) \)
   - Right half using lower bound \( opt(i, n/2) \)

By maintaining valid bounds on \( opt \), we reduce the total complexity to:

\[
O(m n \log n)
\]

Each possible value of \( opt(i, j) \) appears in at most \( \log n \) recursion levels.

Importantly, the method does not depend on how balanced the splits are.  
Across each recursion level, every candidate \( k \) is used at most twice.

\subsection*{Generic Implementation}

Below is a common template.  
The function \texttt{compute} calculates one row \( i \) of DP values (\texttt{dp\_cur})  
using the previous row (\texttt{dp\_before}).  

It must be called as:

\[
\texttt{compute(0, n-1, 0, n-1)}
\]

The function \texttt{solve} computes all \( m \) rows.

\begin{lstlisting}
int m, n;
vector<long long> dp_before, dp_cur;

long long C(int i, int j);

// compute dp_cur[l] ... dp_cur[r]
void compute(int l, int r, int optl, int optr) {
    if (l > r)
        return;

    int mid = (l + r) >> 1;
    pair<long long, int> best = {LLONG_MAX, -1};

    for (int k = optl; k <= min(mid, optr); k++) {
        best = min(best, {
            (k ? dp_before[k - 1] : 0) + C(k, mid),
            k
        });
    }

    dp_cur[mid] = best.first;
    int opt = best.second;

    compute(l, mid - 1, optl, opt);
    compute(mid + 1, r, opt, optr);
}

long long solve() {
    dp_before.assign(n, 0);
    dp_cur.assign(n, 0);

    for (int i = 0; i < n; i++)
        dp_before[i] = C(0, i);

    for (int i = 1; i < m; i++) {
        compute(0, n - 1, 0, n - 1);
        dp_before = dp_cur;
    }

    return dp_before[n - 1];
}
\end{lstlisting}

\subsection*{Things to Watch Out For}

The hardest part of Divide and Conquer DP is proving the monotonicity of \( opt \).

A common sufficient condition is the \textbf{quadrangle inequality}:

\[
C(a, c) + C(b, d) \le C(a, d) + C(b, c)
\quad \text{for all } a \le b \le c \le d
\]

Many Divide and Conquer DP problems can also be solved using the  
\textbf{Convex Hull Trick}, and vice versa.  
Understanding both techniques is highly recommended.
\hline
\section*{Möbius Inversion}
\hline
\subsection*{Möbius Function}

The Möbius function $\mu(n)$ is defined as
\[
\mu(n)=
\begin{cases}
1 & \text{if } n=1, \\
(-1)^k & \text{if } n \text{ is a product of } k \text{ distinct primes}, \\
0 & \text{if } n \text{ has a squared prime factor}.
\end{cases}
\]

Key identity:
\[
\sum_{d \mid n} \mu(d) =
\begin{cases}
1 & \text{if } n=1, \\
0 & \text{if } n>1.
\end{cases}
\]
\subsection*{Inversion Formula}

If
\[
F(n) = \sum_{d \mid n} f(d),
\]
then
\[
f(n) = \sum_{d \mid n} \mu(d)\,F\!\left(\frac{n}{d}\right).
\]

Equivalently, in Dirichlet convolution form:
\[
F = f * 1
\quad \Longleftrightarrow \quad
f = F * \mu.
\]

\subsection*{Common CP Applications}

Counting pairs with $\gcd(i,j)=1$:
\[
\sum_{d=1}^{n}
\mu(d)
\left\lfloor \frac{n}{d} \right\rfloor^2
\]

Used to:
\begin{itemize}
    \item Remove overcounting over divisors
    \item Transform sums over $\gcd$
    \item Invert divisor summations
\end{itemize}
\subsection*{Linear Sieve for $\mu$ (O(n))}

\begin{lstlisting}
const int N = 1e6;
int mu[N], primes[N], pcnt;
bool is_comp[N];

void mobius_sieve(int n) {
    mu[1] = 1;
    for (int i = 2; i <= n; i++) {
        if (!is_comp[i]) {
            primes[pcnt++] = i;
            mu[i] = -1;
        }
        for (int j = 0; j < pcnt && i * primes[j] <= n; j++) {
            int x = i * primes[j];
            is_comp[x] = true;
            if (i % primes[j] == 0) {
                mu[x] = 0;
                break;
            } else {
                mu[x] = -mu[i];
            }
        }
    }
}
\end{lstlisting}
\hline
\section{RAND mt19937}
\hline
\begin{lstlisting}
#include <bits/stdc++.h>
using namespace std;

mt19937 rng(chrono::steady_clock::now().time_since_epoch().count());

int main() {
    int x = rng();               // random 32-bit integer
    int y = rng() % 100;         // 0 to 99
    cout << x << " " << y << "\n";
}
\end{lstlisting}
\hline
\section{RAND mt19937 Faster}
\hline
\begin{lstlisting}
mt19937_64 rng64(chrono::steady_clock::now().time_since_epoch().count());
unsigned long long x = rng64();
\end{lstlisting}
\hline
\section{GNU Extensions}
\hline
\subsection*{Ordered Set / Ordered Map (PBDS)}

Header:
\begin{lstlisting}
#include <ext/pb_ds/assoc_container.hpp>
#include <ext/pb_ds/tree_policy.hpp>
using namespace __gnu_pbds;
\end{lstlisting}

\textbf{Ordered Set}

\begin{lstlisting}
template<typename T>
using ordered_set = tree<
    T,
    null_type,
    less<T>,
    rb_tree_tag,
    tree_order_statistics_node_update>;
\end{lstlisting}

\textbf{Ordered Map}

\begin{lstlisting}
template<typename K, typename V>
using ordered_map = tree<
    K,
    V,
    less<K>,
    rb_tree_tag,
    tree_order_statistics_node_update>;
\end{lstlisting}

\textbf{Supported operations}

\begin{lstlisting}
ordered_set<int> s;

s.insert(x);
s.erase(x);

// number of elements strictly smaller than x
int k = s.order_of_key(x);

// iterator to k-th (0-indexed) element
auto it = s.find_by_order(k);
\end{lstlisting}

Time complexity: \(O(\log n)\)

---

\subsection*{Ordered Multiset}

PBDS does not support duplicates directly.  
Trick: store pairs.

\begin{lstlisting}
using pii = pair<int,int>;

ordered_set<pii> ms;
int id = 0;

ms.insert({value, id++});

// erase one occurrence
ms.erase(ms.lower_bound({value, 0}));
\end{lstlisting}

Alternative trick (less clean):
\begin{lstlisting}
tree<int, null_type, less_equal<int>,
     rb_tree_tag,
     tree_order_statistics_node_update> ms;
\end{lstlisting}

---

\subsection*{Rope}

Header:
\begin{lstlisting}
#include <ext/rope>
using namespace __gnu_cxx;
\end{lstlisting}

Rope is useful for fast string modifications.

\begin{lstlisting}
rope<char> r;

r.insert(pos, "abc");
r.erase(pos, len);
r.replace(pos, len, "xyz");

char c = r[i];
\end{lstlisting}

Time complexity:
\[
\text{insert / erase / substring } = O(\log n)
\]

Common use cases:
\begin{itemize}
    \item Editor simulation
    \item Large string modifications
    \item Persistent-like behavior
\end{itemize}

\hline
\section{Debug}
\hline
\begin{verbatim}

Pre-submit:
Write a few simple test cases if sample is not enough.
Are time limits close? If so, generate max cases.
Is the memory usage fine?
Could anything overflow?
Make sure to submit the right file.

Wrong answer:
Print your solution! Print debug output, as well.
Are you clearing all data structures between test cases?
Can your algorithm handle the whole range of input?
Read the full problem statement again.
Do you handle all corner cases correctly?
Have you understood the problem correctly?
Any uninitialized variables?
Any overflows?
Confusing N and M, i and j, etc.?
Are you sure your algorithm works?
What special cases have you not thought of?
Are you sure the STL functions you use work as you think?
Add some assertions, maybe resubmit.
Create some testcases to run your algorithm on.
Go through the algorithm for a simple case.
Go through this list again.
Explain your algorithm to a teammate.
Ask the teammate to look at your code.
Go for a small walk, e.g. to the toilet.
Is your output format correct? (including whitespace)
Rewrite your solution from the start or let a teammate do it.

Runtime error:
Have you tested all corner cases locally?
Any uninitialized variables?
Are you reading or writing outside the range of any vector?
Any assertions that might fail?
Any possible division by 0? (mod 0 for example)
Any possible infinite recursion?
Invalidated pointers or iterators?
Are you using too much memory?
Debug with resubmits (e.g. remapped signals, see Various).

Time limit exceeded:
Do you have any possible infinite loops?
What is the complexity of your algorithm?
Are you copying a lot of unnecessary data? (References)
How big is the input and output? (consider scanf)
Avoid vector, map. (use arrays/unordered_map)
What do your teammates think about your algorithm?

Memory limit exceeded:
What is the max amount of memory your algorithm should need?
Are you clearing all data structures between test cases?
\end{verbatim}

\section{Zeta and Morbius}
\begin{lstlisting}
template<typename T>
void SubsetZetaTransform(vector<T>& v) {
	const int n = v.size(); // n must be a power of 2
	for (int j = 1; j < n; j <<= 1) {
		for (int i = 0; i < n; i++)
			if (i & j) v[i] += v[i ^ j];
	}
}

template<typename T>
void SubsetMobiusTransform(vector<T>& v) {
	const int n = v.size(); // n must be a power of 2
	for (int j = 1; j < n; j <<= 1) {
		for (int i = 0; i < n; i++)
			if (i & j) v[i] -= v[i ^ j];
	}
}

template<typename T>
void SupersetZetaTransform(vector<T>& v) {
	const int n = v.size(); // n must be a power of 2
	for (int j = 1; j < n; j <<= 1) {
		for (int i = 0; i < n; i++)
			if (i & j) v[i ^ j] += v[i];
	}
}

template<typename T>
void SupersetMobiusTransform(vector<T>& v) {
	const int n = v.size(); // n must be a power of 2
	for (int j = 1; j < n; j <<= 1) {
		for (int i = 0; i < n; i++)
			if (i & j) v[i ^ j] -= v[i];
	}
}

template<typename T>
void DivisorZetaTransform(vector<T>& v) {
	const int n = (int)v.size() - 1;
	for (int p : PrimeEnumerate(n)) {
		for (int i = 1; i * p <= n; i++)
			v[i * p] += v[i];
	}
}

template<typename T>
void DivisorMobiusTransform(vector<T>& v) {
	const int n = (int)v.size() - 1;
	for (int p : PrimeEnumerate(n)) {
		for (int i = n / p; i; i--)
			v[i * p] -= v[i];
	}
}
template<typename T>
void MultipleZetaTransform(vector<T>& v) {
	const int n = (int)v.size() - 1;
	for (int p : PrimeEnumerate(n)) {
		for (int i = n / p; i; i--)
			v[i] += v[i * p];
	}
}

template<typename T>
void MultipleMobiusTransform(vector<T>& v) {
	const int n = (int)v.size() - 1;
	for (int p : PrimeEnumerate(n)) {
		for (int i = 1; i * p <= n; i++)
			v[i] -= v[i * p];
	}

\end{lstlisting}

\section{Hadamard Transform}
\begin{lstlisting}
poly FWHT(poly P, bool inverse) {
    for (len = 1; 2 * len <= degree(P); len <<= 1) {
        for (i = 0; i < degree(P); i += 2 * len) {
            for (j = 0; j < len; j++) {
                u = P[i + j];
                v = P[i + len + j];
                P[i + j] = u + v;
                P[i + len + j] = u - v;
            }
        }
    }
    if (inverse) {
        for (i = 0; i < degree(P); i++)
            P[i] = P[i] / degree(P);
    }
    return P;
}
\end{lstlisting}

\section{Priority Queue with Delete}
\begin{lstlisting}
struct Heap {
	PQ<int> qi,qo;
	void ins(int v,int o) { (o>0?qi:qo).push(v); }
	int top() {
		while(qo.size()&&qi.top()==qo.top()) qi.pop(),qo.pop();
		return qi.top();
	}
};
\end{lstlisting}

\section{DSU on Segtree}
\begin{lstlisting}
struct node {
	int s, l, r;
	node() : s(0), l(0), r(0) {}
} tree[40 * maxn];
int tp = 1;
void upd(int &u, int tl, int tr, int pos, int delta) 
{
	if(!u) u = tp++;
	tree[u].s += delta;
	if(tl < tr) {
		int tm = (tl + tr) / 2;
		if(pos <= tm) upd(tree[u].l, tl, tm, pos, delta);
		else upd(tree[u].r, tm + 1, tr, pos, delta);
	}
}
int query(int u, int tl, int tr, int l, int r)
{
	if(l > r || u == 0) return 0;
	if(l == tl && r == tr) return tree[u].s;
	int tm = (tl + tr) / 2;
	return query(tree[u].l, tl, tm, l, min(r, tm)) +
			query(tree[u].r, tm + 1, tr, max(l, tm + 1), r);
}
int merge(int u, int v, int tl, int tr)
{
	if(!u || !v) return u | v;
	tree[u].s += tree[v].s;
	if(tl < tr) {
		int tm = (tl + tr) / 2;
		tree[u].l = merge(tree[u].l, tree[v].l, tl, tm);
		tree[u].r = merge(tree[u].r, tree[v].r, tm + 1, tr);
	}
	return u;
}
\end{lstlisting}

\section{Treap}
\begin{lstlisting}
typedef struct item * pitem;
struct item {
    int prior, value, cnt;
    bool rev;
    pitem l, r;
};

int cnt (pitem it) {
    return it ? it->cnt : 0;
}

void upd_cnt (pitem it) {
    if (it)
        it->cnt = cnt(it->l) + cnt(it->r) + 1;
}

void push (pitem it) {
    if (it && it->rev) {
        it->rev = false;
        swap (it->l, it->r);
        if (it->l)  it->l->rev ^= true;
        if (it->r)  it->r->rev ^= true;
    }
}

void merge (pitem & t, pitem l, pitem r) {
    push (l);
    push (r);
    if (!l || !r)
        t = l ? l : r;
    else if (l->prior > r->prior)
        merge (l->r, l->r, r),  t = l;
    else
        merge (r->l, l, r->l),  t = r;
    upd_cnt (t);
}

void split (pitem t, pitem & l, pitem & r, int key, int add = 0) {
    if (!t)
        return void( l = r = 0 );
    push (t);
    int cur_key = add + cnt(t->l);
    if (key <= cur_key)
        split (t->l, l, t->l, key, add),  r = t;
    else
        split (t->r, t->r, r, key, add + 1 + cnt(t->l)),  l = t;
    upd_cnt (t);
}

void reverse (pitem t, int l, int r) {
    pitem t1, t2, t3;
    split (t, t1, t2, l);
    split (t2, t2, t3, r-l+1);
    t2->rev ^= true;
    merge (t, t1, t2);
    merge (t, t, t3);
}

void output (pitem t) {
    if (!t)  return;
    push (t);
    output (t->l);
    printf ("%d ", t->value);
    output (t->r);
}
\end{lstlisting}
\end{document}